% !PRINTED 325
\weberpage{325}
%<<CONT
and consequently formula (12) is proved if it holds for any two prime discriminants.

But by (10), if $p$ and $p'$ are odd primes:
\[
(\pm p, \pm p') = \pm\left(\frac{p'}{p}\right), \quad (\pm p', \pm p) = \pm\left(\frac{p'}{p}\right), % [sic] The source prints p'/p in both members; the following prose requires p/p' in the second.
\]
in which the negative sign stands on the right only when \gesperrt{both signs on the left are negative}, that is, when $p$ as well as $p' \equiv 3 \pmod{4}$. In this case, and only in this case, $\left(\dfrac{p'}{p}\right) = -\left(\dfrac{p}{p'}\right)$, otherwise $\left(\dfrac{p'}{p}\right) = \left(\dfrac{p}{p'}\right)$. Thus formula (13) proves to be correct for this case.

Further, by (11) and (3), (5)
\[
(-4, \pm p) = +1, \quad (\pm p, -4) = \pm 1,
\]
\[
(p, \pm 8) = (p, \pm 2) = (p, 2) = (-1)^{\frac{p^2-1}{8}},
\]
\[
(\pm 8, p) = \left(\frac{\pm 8}{p}\right) = (-1)^{\frac{p^2-1}{8}},
\]
\[
[p \equiv 1 \pmod{4}],
\]
\[
(-p, \pm 8) = (-p, \pm 2) = \pm(-p, 2) = -1,
\]
\[
(\pm 8, -p) = \pm\left(\frac{\pm 8}{p}\right) = -1,
\]
\[
[p \equiv 3 \pmod{8}],
\]
from which formula (13) follows for these cases, and it is thereby proved in general.

From (13) it follows further, since every square number is a positive discriminant, that if $m$ is \gesperrt{relatively prime} to $D$,
\weq{14}{(D, m^2) = 1}
and consequently, by (8):
\weq{15}{(D, m^2 n) = (D, n),}
if $m$ and $D$ have no common divisor.

We have further the theorem:
\weq{16}{(D, n) = (D, n'),}
if $n \equiv n' \pmod{D}$.

The theorem is correct when $n$ and $D$ have a common divisor, and by (12) it is proved in general if it holds for every prime discriminant $D$. For these, however, it follows at once from (10), (11).
